\makeatletter
\def\input@path{{./}{../../../preamble/}{book2026-2027/preamble/}}
\makeatother
\input{preamble.tex}
\input{graphs/highlights.tex}
\input{graphs/G10_T2_bar_graphs.tex}
\usepackage{needspace}
\hypersetup{pdftitle={Grade 10 Term 2 Discrete Probability}}
\setlength{\emergencystretch}{2em}
\title{Grade 10, Term 2, Lesson 1, Criterion C1\\Outcomes, Support and Probability Mass Functions}
\author{}
\date{}
\begin{document}
\maketitle
\setcounter{tocdepth}{1}
\setcounter{secnumdepth}{2}
\phantomsection\label{toc}
\tableofcontents
\section*{Evaluated criterion}
\textbf{C1:} Characterises discrete random variables in economic and financial contexts by identifying their support, probability mass function and possible outcomes.
\section{Introduction}\label{sec:t2-c1-introduction}
All data in this activity are hypothetical teaching models. Probabilities have no monetary units. Keep exact expressions where useful and round final probabilities to four decimal places.

\GTenHighlightLegend

\subsection{From a school sale to a number}

A school enterprise offers a notebook to each of three visitors. Each visitor either buys (B) or does not buy (N). An \emph{outcome} records the whole story, such as BNN: the first visitor buys, while the second and third do not. A \emph{random variable} assigns a number to every outcome.

The experiment has eight possible outcomes:
\[
\Omega=\{\mathrm{NNN},\mathrm{BNN},\mathrm{NBN},\mathrm{NNB},
\mathrm{BBN},\mathrm{BNB},\mathrm{NBB},\mathrm{BBB}\}.
\]

Consider three random variables defined on this experiment:

\begin{itemize}
	\item \textbf{Example 1 (Number of buyers).}
	Let $\GTenHLValue{X}$ be the \GTenHLValue{number of visitors who buy a notebook}.
	Then $\GTenHLValue{X}(\GTenHLOutcome{\mathrm{BNN}})=\GTenHLValue{1}$, $\GTenHLValue{X}(\GTenHLOutcome{\mathrm{BBN}})=\GTenHLValue{2}$,
	$\GTenHLValue{X}(\GTenHLOutcome{\mathrm{BBB}})=\GTenHLValue{3}$, and $\GTenHLValue{X}(\GTenHLOutcome{\mathrm{NNN}})=\GTenHLValue{0}$.
	
	\item \textbf{Example 2 (Sales revenue).}
	If each notebook costs \$5, let $\GTenHLSecond{Y}$ be the \GTenHLSecond{total revenue}.
	Then $\GTenHLSecond{Y}(\GTenHLOutcome{\mathrm{BNN}})=\GTenHLSecond{5}$, $\GTenHLSecond{Y}(\GTenHLOutcome{\mathrm{BBN}})=\GTenHLSecond{10}$,
	$\GTenHLSecond{Y}(\GTenHLOutcome{\mathrm{BBB}})=\GTenHLSecond{15}$, and $\GTenHLSecond{Y}(\GTenHLOutcome{\mathrm{NNN}})=\GTenHLSecond{0}$.
	
	\item \textbf{Example 3 (First visitor's purchase).}
	Let $\GTenHLThird{Z}$ equal \GTenHLThird{1} if the \GTenHLThird{first visitor buys a notebook}
	and \GTenHLThird{0} otherwise. Then $\GTenHLThird{Z}(\GTenHLOutcome{\mathrm{BNN}})=\GTenHLThird{1}$,
	$\GTenHLThird{Z}(\GTenHLOutcome{\mathrm{BBN}})=\GTenHLThird{1}$, $\GTenHLThird{Z}(\GTenHLOutcome{\mathrm{BBB}})=\GTenHLThird{1}$,
	and $\GTenHLThird{Z}(\GTenHLOutcome{\mathrm{NBN}})=\GTenHLThird{0}$.
\end{itemize}

The notation $\GTenHLValue{X}(\GTenHLOutcome{\mathrm{BNN}})=\GTenHLValue{1}$ is read ``X of BNN equals one''. The capital letter names the random variable (the rule), while a lowercase letter such as $x$ represents a possible or observed numerical value.

\textbf{Check.} Are BNN and NBN the same outcome? Do they produce the same values of $\GTenHLValue{X}$, $\GTenHLSecond{Y}$, and $\GTenHLThird{Z}$?

\emph{Suggested answer:} No; they describe different visitors buying. Both give $\GTenHLValue{X}=\GTenHLValue{1}$ and $\GTenHLSecond{Y}=\GTenHLSecond{5}$, but $\GTenHLThird{Z}(\GTenHLOutcome{\mathrm{BNN}})=\GTenHLThird{1}$ whereas $\GTenHLThird{Z}(\GTenHLOutcome{\mathrm{NBN}})=\GTenHLThird{0}$.


\subsection{Support, measurements and units}

The \emph{support} of a random variable is the set of values with positive probability. Different random variables describing the same experiment may have different supports and measurement units.

Consider the three random variables introduced previously:

\begin{itemize}
	\item \textbf{Example 1 (Number of buyers).}
	The support of $\GTenHLValue{X}$ is $\GTenHLValue{\{0,1,2,3\}}$, measured in \GTenHLValue{buyers}.
	The values represent zero, one, two, or three visitors
	purchasing a notebook.
	
	\item \textbf{Example 2 (Sales revenue).}
	The support of $\GTenHLSecond{Y}$ is $\GTenHLSecond{\{0,5,10,15\}}$, measured in \GTenHLSecond{dollars}.
	Revenue increases in increments of \$5 because each
	notebook has a fixed price.
	
	\item \textbf{Example 3 (First visitor's purchase).}
	The support of $\GTenHLThird{Z}$ is $\GTenHLThird{\{0,1\}}$, with no physical unit.
	The values indicate whether the first visitor buys (1)
	or does not buy (0).
\end{itemize}

Braces denote a set of possible values. All three examples have \emph{discrete} supports because their values can be listed individually. A discrete support may be finite, such as $\{0,1,2,3\}$, or countably infinite, such as $\{1,2,3,\ldots\}$, where the dots indicate that the pattern continues indefinitely.

A \emph{continuous} random variable, by contrast, can take values throughout an interval. For example, the exact time required to complete a purchase may be modelled as continuous and measured in seconds. If times are rounded to whole seconds, the recorded values become discrete.

\textbf{Check.} Can $\GTenHLValue{X}=\GTenHLValue{1.6}$, $\GTenHLSecond{Y}=\GTenHLSecond{7}$, or $\GTenHLThird{Z}=\GTenHLThird{2}$? Explain.

\emph{Suggested answer:} No. The buyer count $\GTenHLValue{X}$ must be an integer between 0 and 3; revenue $\GTenHLSecond{Y}$ can only be 0, 5, 10, or 15 dollars; and the indicator $\GTenHLThird{Z}$ can only be 0 or 1.


\subsection{Assign probability to each value}

The \emph{probability mass function} (PMF) assigns a probability to each value of a discrete random variable. For $\GTenHLValue{X}$, it is written $p_X(x)=P(X=x)$, read ``the probability that X equals x''. The subscript identifies the random variable: $p_Y(y)=P(Y=y)$ describes revenue, while $p_Z(z)=P(Z=z)$ describes the first visitor's purchase indicator.

Every probability must be non-negative, and the probabilities over the support must add to 1. A value outside the support has probability zero. To find the probability of a value, add the probabilities of all outcomes that produce it; different outcomes need not be equally likely.

For each visitor in the school sale, assign
\[
P(\GTenHLOutcome{\mathrm{B}})=0.4,\qquad P(\GTenHLOutcome{\mathrm{N}})=0.6,
\]
where B means buying a notebook and N means not buying. These probabilities add to 1. Assume that the three visitors decide independently and that each has these same probabilities. Independence means that one visitor's decision does not change the probabilities for another visitor.

\textbf{Step 1: Find the probability of each complete outcome.} Multiply the probabilities of the three decisions in their recorded order. For example,
\[
\begin{aligned}
P(\GTenHLOutcome{\mathrm{BNN}})&=0.4\times0.6\times0.6=\GTenHLProbability{0.144},\\
P(\GTenHLOutcome{\mathrm{BBN}})&=0.4\times0.4\times0.6=\GTenHLProbability{0.096}.
\end{aligned}
\]
Applying the same rule to all eight outcomes gives
\[
\begin{aligned}
	P(\GTenHLOutcome{\mathrm{NNN}})&=0.6^3=\GTenHLProbability{0.216}, & P(\GTenHLOutcome{\mathrm{BNN}})&=0.4\times0.6^2=\GTenHLProbability{0.144},\\
	P(\GTenHLOutcome{\mathrm{NBN}})&=0.6\times0.4\times0.6=\GTenHLProbability{0.144}, & P(\GTenHLOutcome{\mathrm{NNB}})&=0.6^2\times0.4=\GTenHLProbability{0.144},\\
	P(\GTenHLOutcome{\mathrm{BBN}})&=0.4^2\times0.6=\GTenHLProbability{0.096}, & P(\GTenHLOutcome{\mathrm{BNB}})&=0.4\times0.6\times0.4=\GTenHLProbability{0.096},\\
	P(\GTenHLOutcome{\mathrm{NBB}})&=0.6\times0.4^2=\GTenHLProbability{0.096}, & P(\GTenHLOutcome{\mathrm{BBB}})&=0.4^3=\GTenHLProbability{0.064}.
\end{aligned}
\]
The outcomes are not all equally likely: their probabilities come from the independent decisions, not from dividing 1 by the number of outcomes. Check the total:
\[
\GTenHLProbability{0.216}+3(0.144)+3(0.096)+0.064=1.
\]

\textbf{Step 2: Group outcomes by the value of each variable.} Add the probabilities of the outcomes that give the same value.

\begin{itemize}
	\item \textbf{Example 1 (Number of buyers).}
	The outcomes BNN, NBN, and NNB produce $\GTenHLValue{X}=\GTenHLValue{1}$, so
	\[
	\begin{aligned}
		p_{\GTenHLValue{X}}(\GTenHLValue{1})
		&=P(\GTenHLOutcome{\mathrm{BNN}})+P(\GTenHLOutcome{\mathrm{NBN}})+P(\GTenHLOutcome{\mathrm{NNB}})\\
		&=\GTenHLProbability{0.144}+\GTenHLProbability{0.144}+\GTenHLProbability{0.144}=\GTenHLProbability{0.432}.
	\end{aligned}
	\]
	Similarly, BBN, BNB, and NBB produce $\GTenHLValue{X}=\GTenHLValue{2}$, so
	\[
	p_{\GTenHLValue{X}}(\GTenHLValue{2})=\GTenHLProbability{0.096}+\GTenHLProbability{0.096}+\GTenHLProbability{0.096}=\GTenHLProbability{0.288}.
	\]
	Only NNN gives $\GTenHLValue{X}=\GTenHLValue{0}$, and only BBB gives $\GTenHLValue{X}=\GTenHLValue{3}$.
	The complete PMF is
	\[
	\begin{aligned}
		p_{\GTenHLValue{X}}(\GTenHLValue{0})&=\GTenHLProbability{0.216}, & p_{\GTenHLValue{X}}(\GTenHLValue{1})&=\GTenHLProbability{0.432},\\
		p_{\GTenHLValue{X}}(\GTenHLValue{2})&=\GTenHLProbability{0.288}, & p_{\GTenHLValue{X}}(\GTenHLValue{3})&=\GTenHLProbability{0.064}.
	\end{aligned}
	\]
	
	\item \textbf{Example 2 (Sales revenue).}
	Each notebook sells for 5 USD, so $\GTenHLSecond{Y}=5\GTenHLValue{X}$ when revenue
	is recorded in USD. The same outcomes that produce one
	buyer produce revenue of 5 USD. Thus the complete PMF is
	\[
	\begin{aligned}
		p_{\GTenHLSecond{Y}}(\GTenHLSecond{0})&=\GTenHLProbability{0.216}, & p_{\GTenHLSecond{Y}}(\GTenHLSecond{5})&=\GTenHLProbability{0.432},\\
		p_{\GTenHLSecond{Y}}(\GTenHLSecond{10})&=\GTenHLProbability{0.288}, & p_{\GTenHLSecond{Y}}(\GTenHLSecond{15})&=\GTenHLProbability{0.064}.
	\end{aligned}
	\]
	The revenue values have monetary units; their
	probabilities do not.
	
	\item \textbf{Example 3 (First visitor's purchase).}
	The first visitor buys in BNN, BBN, BNB, and BBB, so
	\[
	\begin{aligned}
		p_{\GTenHLThird{Z}}(\GTenHLThird{1})
		&=\GTenHLProbability{0.144}+\GTenHLProbability{0.096}+\GTenHLProbability{0.096}+\GTenHLProbability{0.064}\\
		&=\GTenHLProbability{0.400}.
	\end{aligned}
	\]
	The first visitor does not buy in NNN, NBN, NNB,
	and NBB, so
	\[
	\begin{aligned}
		p_{\GTenHLThird{Z}}(\GTenHLThird{0})
		&=\GTenHLProbability{0.216}+\GTenHLProbability{0.144}+\GTenHLProbability{0.144}+\GTenHLProbability{0.096}\\
		&=\GTenHLProbability{0.600}.
	\end{aligned}
	\]
\end{itemize}

\textbf{Step 3: Check and interpret.} For $\GTenHLValue{X}$ and $\GTenHLSecond{Y}$, the probabilities total $\GTenHLProbability{0.216}+\GTenHLProbability{0.432}+\GTenHLProbability{0.288}+\GTenHLProbability{0.064}=1$; for $\GTenHLThird{Z}$, they total $\GTenHLProbability{0.600}+\GTenHLProbability{0.400}=1$. All entries are non-negative. The probabilities for $\GTenHLThird{Z}$ agree with the original chances of the first visitor not buying or buying. Although $\GTenHLValue{X}$ and $\GTenHLSecond{Y}$ have different supports and units, corresponding values have the same probabilities because $\GTenHLSecond{Y}=5\GTenHLValue{X}$. The variable $\GTenHLThird{Z}$ groups the outcomes differently.


% Graphs follow completed calculations; captions stay with their panels.
\par\noindent \begin{minipage}[t]{.485\linewidth}
\centering
\GTenOutcomeProbabilityGraph[width=\linewidth,height=4.7cm,rotate=35,ystep=.05]{NNN,BNN,NBN,NNB,BBN,BNB,NBB,BBB}{.216,.144,.144,.144,.096,.096,.096,.064}
\par\small \textbf{School sale: outcome probabilities.} Each bar describes one complete three-visitor outcome.
\end{minipage}\hfill\begin{minipage}[t]{.485\linewidth}
\centering
\GTenOutcomeValueGraph[width=\linewidth,height=4.7cm,rotate=35,ylabel={$X(\omega)$ (buyers)}]{NNN,BNN,NBN,NNB,BBN,BNB,NBB,BBB}{0,1,1,1,2,2,2,3}
\par\small \textbf{Assigned buyer counts.} The height is $X(\omega)$ in buyers, not a probability.
\end{minipage}\par\medskip
\par\noindent \begin{minipage}[t]{.485\linewidth}
\centering
\GTenOutcomeValueGraph[width=\linewidth,height=4.7cm,rotate=35,ystep=5,ylabel={$Y(\omega)$ (USD)}]{NNN,BNN,NBN,NNB,BBN,BNB,NBB,BBB}{0,5,5,5,10,10,10,15}
\par\small \textbf{Assigned revenue.} The height is $Y(\omega)=5X(\omega)$ in USD.
\end{minipage}\hfill\begin{minipage}[t]{.485\linewidth}
\centering
\GTenOutcomeValueGraph[width=\linewidth,height=4.7cm,rotate=35,ylabel={$Z(\omega)$ (indicator)}]{NNN,BNN,NBN,NNB,BBN,BNB,NBB,BBB}{0,1,0,0,1,1,0,1}
\par\small \textbf{Assigned indicator.} $Z(\omega)$ is 1 exactly when the first visitor buys; it has no units.
\end{minipage}\par\medskip
\par\noindent \begin{minipage}[t]{.32\linewidth}
\centering
\GTenPMFGraph[width=\linewidth,height=4.7cm,xlabel={$x$ (buyers)}]{0,1,2,3}{.216,.432,.288,.064}
\par\small \textbf{Buyer-count PMF.} Group outcomes by $\GTenHLValue{X}$ and add their probabilities.
\end{minipage}\hfill\begin{minipage}[t]{.32\linewidth}
\centering
\GTenPMFGraph[width=\linewidth,height=4.7cm,xlabel={$y$ (USD)},xstep=5]{0,5,10,15}{.216,.432,.288,.064}
\par\small \textbf{Revenue PMF.} Group outcomes by $\GTenHLSecond{Y}$; probabilities have no USD units.
\end{minipage}\hfill\begin{minipage}[t]{.32\linewidth}
\centering
\GTenPMFGraph[width=\linewidth,height=4.7cm,xlabel={$z$ (indicator)}]{0,1}{.6,.4}
\par\small \textbf{First-visitor PMF.} Group outcomes by $\GTenHLThird{Z}$; these bars give probabilities of indicator values.
\end{minipage}\par\medskip


\textbf{Check.} Is $p_{\GTenHLSecond{Y}}(\GTenHLSecond{5})=\GTenHLProbability{0.432}$ a payment of \GTenHLProbability{0.432} USD? What does it mean?

\emph{Suggested answer:} No. It is a dimensionless probability: there is a \GTenHLProbability{43.2\%} chance that revenue equals \GTenHLSecond{5 USD}.

\textbf{Check.} Are the events $\GTenHLValue{X}=\GTenHLValue{1}$, $\GTenHLSecond{Y}=\GTenHLSecond{5}$, and $\GTenHLThird{Z}=\GTenHLThird{1}$ the same?

\emph{Suggested answer:} The first two are the same event, consisting of BNN, NBN, and NNB. The event $\GTenHLThird{Z}=\GTenHLThird{1}$ consists of BNN, BBN, BNB, and BBB, so it is different. Their probabilities are \GTenHLProbability{0.432}, \GTenHLProbability{0.432}, and \GTenHLProbability{0.400}, respectively.


\Needspace{10\baselineskip}
\section{General worked problems}
\Needspace{8\baselineskip}
\subsection{Introductory general problem: grouping outcomes}
\subsubsection*{Mathematical wording}
An experiment has outcomes $\GTenHLOutcome{a},\GTenHLOutcome{b},\GTenHLOutcome{c},\GTenHLOutcome{d}$ with probabilities $\GTenHLProbability{0.1},\GTenHLProbability{0.2},\GTenHLProbability{0.3},\GTenHLProbability{0.4}$. Define $\GTenHLValue{X}(\GTenHLOutcome{a})=\GTenHLValue{0}$, $\GTenHLValue{X}(\GTenHLOutcome{b})=\GTenHLValue{X}(\GTenHLOutcome{c})=\GTenHLValue{1}$, $\GTenHLValue{X}(\GTenHLOutcome{d})=\GTenHLValue{2}$.

\textbf{Questions/Tasks}
\begin{enumerate}
\item Construct the mapping and PMF; state the support and find $P(X=1)$.
\end{enumerate}
\Needspace{6\baselineskip}
\subsubsection*{Financial narrative wording}
A school shop receives one payment. There are exactly four possible payment methods: \GTenHLOutcome{cash}, \GTenHLOutcome{debit card}, \GTenHLOutcome{credit card}, or \GTenHLOutcome{a mobile wallet}. The chance of cash is \GTenHLProbability{10\%}, of debit card is \GTenHLProbability{20\%}, of credit card is \GTenHLProbability{30\%}, and of a mobile wallet is \GTenHLProbability{40\%}. The shop pays no processing fee for cash, a fee of \GTenHLValue{\$1} for either type of card, and a fee of \GTenHLValue{\$2} for a mobile wallet. These are the only processing charges.

\textbf{Questions/Tasks}
\begin{enumerate}
\item Show which processing fee goes with each payment method. List all possible fee amounts and give the chance of each amount. What is the chance that the shop pays exactly \$1 in processing fees?
\end{enumerate}
\subsubsection*{Complete worked solution}
\begin{center}
\begin{tabular}{llp{0.48\linewidth}}
\toprule
Symbol & Value & Meaning in this problem \\ \midrule
$\GTenHLValue{X}$ & $0,1,2$ & Possible numerical values; fees in USD in the story \\
$p_X(1)$ & $0.2+0.3$ & Mass grouped from two outcomes \\
\bottomrule
\end{tabular}
\end{center}

\par\noindent\begin{minipage}{\linewidth}
\textbf{Step 1: Define the variable.} $\GTenHLValue{X}$ is the numerical score assigned to an outcome; its units are score units in the abstract version. In the financial story, define $\GTenHLValue{X}$ as the \GTenHLValue{processing fee in USD}. Label cash, debit card, credit card, and mobile wallet by $\GTenHLOutcome{a},\GTenHLOutcome{b},\GTenHLOutcome{c},\GTenHLOutcome{d}$, respectively; their probabilities are $\GTenHLProbability{0.1},\GTenHLProbability{0.2},\GTenHLProbability{0.3},\GTenHLProbability{0.4}$. The fees give exactly the same numerical mapping.
\end{minipage}\par\medskip
\par\noindent\begin{minipage}{\linewidth}
\textbf{Step 2: Map outcomes.} \begin{center}
\begin{tabular}{lllll}
\toprule
Outcome & $\GTenHLOutcome{a}$ & $\GTenHLOutcome{b}$ & $\GTenHLOutcome{c}$ & $\GTenHLOutcome{d}$ \\ \midrule
Probability & \GTenHLProbability{0.1} & \GTenHLProbability{0.2} & \GTenHLProbability{0.3} & \GTenHLProbability{0.4} \\
$\GTenHLValue{X}$ value & \GTenHLValue{0} & \GTenHLValue{1} & \GTenHLValue{1} & \GTenHLValue{2} \\
\bottomrule
\end{tabular}
\end{center}
The support is $\GTenHLValue{\{0,1,2\}}$. The outcomes are mutually exclusive and exhaustive; equal likelihood is not assumed.
\end{minipage}\par\medskip
\textbf{Step 3: Translate and calculate.} The rule for grouping outcomes is
\[
p_X(x)=\sum_{\omega:X(\omega)=x}P(\omega).
\]
Read this as: ``The probability that $\GTenHLValue{X}$ takes the value $x$ is the sum of the probabilities of all outcomes that give that value.'' Apply it in the following order:
\begin{enumerate}
\item \textbf{Choose the value $x$.} The capital letter $\GTenHLValue{X}$ names the rule that assigns a number to each outcome; the lowercase letter $x$ is the particular number we want. For example, choose $\GTenHLValue{x}=\GTenHLValue{1}$. In the story, this means a processing fee of \$1.
\item \textbf{Consider each outcome $\omega$.} The Greek letter $\omega$ (omega) stands for one complete outcome, such as $a$, $b$, $c$, or $d$. The expression $X(\omega)$ is the number assigned to that outcome; for example, $\GTenHLValue{X}(\GTenHLOutcome{b})=\GTenHLValue{1}$.
\item \textbf{Keep only outcomes with $X(\omega)=x$.} The colon in $\omega:X(\omega)=x$ means ``such that''. This condition tells us which outcomes to include in the sum. When $\GTenHLValue{x}=\GTenHLValue{1}$, keep $\GTenHLOutcome{b}$ and $\GTenHLOutcome{c}$ because both give 1; leave out $a$ and $d$ because they give 0 and 2.
\item \textbf{Find the probability of each matching outcome.} The term $P(\omega)$ means the probability of the individual outcome, not its assigned numerical value. Here the matching probabilities are $P(\GTenHLOutcome{b})=\GTenHLProbability{0.2}$ and $P(\GTenHLOutcome{c})=\GTenHLProbability{0.3}$.
\item \textbf{Add those probabilities.} The symbol $\sum$ (capital sigma) means ``add''. Different complete outcomes are mutually exclusive, so add their probabilities, counting each matching outcome once:
\[
p_X(\GTenHLValue{1})=P(\GTenHLValue{X}=\GTenHLValue{1})=P(\GTenHLOutcome{b})+P(\GTenHLOutcome{c})=\GTenHLProbability{0.2}+\GTenHLProbability{0.3}=\GTenHLProbability{0.5000}.
\]
The left side names the probability assigned to the value 1; the right side shows how it is built from the outcomes. Financially, add the chances of debit and credit card payments because both produce the same \$1 fee.
\end{enumerate}
Repeat for the other values: $p_X(\GTenHLValue{0})=P(\GTenHLOutcome{a})=\GTenHLProbability{0.1}$ and $p_X(\GTenHLValue{2})=P(\GTenHLOutcome{d})=\GTenHLProbability{0.4}$. If no outcome gives a chosen value, the sum has no terms and its value is zero; for example, $p_X(3)=0$. The complete PMF on the support is:
\begin{center}
\begin{tabular}{llll}
\toprule
$\GTenHLValue{x}$ & \GTenHLValue{0} & \GTenHLValue{1} & \GTenHLValue{2} \\ \midrule
$p_X(x)$ & \GTenHLProbability{0.1} & \GTenHLProbability{0.5} & \GTenHLProbability{0.4} \\
\bottomrule
\end{tabular}
\end{center}

\par\medskip
\par\noindent\begin{minipage}{\linewidth}
\textbf{Step 4: Check and interpret.} $0.1+0.5+0.4=1$ and all entries are non-negative. The score equals 1 in half of repetitions under this model, even though two distinct outcomes produce it. Financially, debit and credit card payments both cost the shop \$1, so the chance of a \$1 processing fee is \GTenHLProbability{50\%}. Fee amounts are measured in USD; probabilities have no monetary units.
\end{minipage}\par\medskip
% Graphs follow completed calculations; captions stay with their panels.
\par\noindent \begin{minipage}[t]{.32\linewidth}
\centering
\GTenOutcomeProbabilityGraph[width=\linewidth,height=4.7cm]{$a$,$b$,$c$,$d$}{.1,.2,.3,.4}
\par\small \textbf{Outcome probabilities.} Distinct outcomes have their own chances.
\end{minipage}\hfill\begin{minipage}[t]{.32\linewidth}
\centering
\GTenOutcomeValueGraph[width=\linewidth,height=4.7cm,ylabel={Assigned score / fee (USD)}]{$a$,$b$,$c$,$d$}{0,1,1,2}
\par\small \textbf{Assigned values.} Both $\GTenHLOutcome{b}$ and $\GTenHLOutcome{c}$ give 1; in the financial wording these are fees in USD.
\end{minipage}\hfill\begin{minipage}[t]{.32\linewidth}
\centering
\GTenPMFGraph[width=\linewidth,height=4.7cm,xlabel={Score / fee (USD)}]{0,1,2}{.1,.5,.4}
\par\small \textbf{Grouped PMF.} The probability at 1 is $0.2+0.3=0.5$.
\end{minipage}\par\medskip


\Needspace{10\baselineskip}
\subsection{Bridge to the general model}
For any discrete variable,
\[p_X(x)=P(X=x)\geq0,\qquad \sum_{x\in S}p_X(x)=1.\]
These statements express two requirements: each assigned probability is non-negative, and the probabilities of all possible values add to 1. Read them step by step:
\begin{enumerate}
\item \textbf{Identify the variable and a value.} The capital letter $X$ names the rule that assigns a number to each outcome; the lowercase letter $x$ represents one possible numerical value. In the preceding grouping problem, $X$ is the score, or the processing fee in USD in the financial story, and its possible values are 0, 1, and 2.
\item \textbf{Read $p_X(x)=P(X=x)$.} The expression $X=x$ describes the event that the variable takes the chosen value. Both $p_X(x)$ and $P(X=x)$ name the probability of that event. For example, $p_X(\GTenHLValue{1})=P(\GTenHLValue{X}=\GTenHLValue{1})=\GTenHLProbability{0.5}$: the chance of a \$1 fee is 50\%. The value has monetary units in the story; its probability has no units.
\item \textbf{Read $\geq0$.} This symbol means ``greater than or equal to zero''. No assigned probability can be negative. In the example, $0.1$, $0.5$, and $0.4$ are all non-negative. Values outside the support have probability zero; for example, $p_X(3)=0$.
\item \textbf{Identify the support $S$.} The support is the set of values with positive probability. Here $S=\GTenHLValue{\{0,1,2\}}$. The notation $x\in S$ means ``x belongs to S'' and tells us which values to use in the sum.
\item \textbf{Expand $\sum_{x\in S}p_X(x)$.} The symbol $\sum$ means ``add''. Let $x$ take each value in $S$ once, find its probability, and add those probabilities. For this support,
% Value-to-mass lookup before the normalization sum.
\[
p_X(\GTenHLValue{0})=\GTenHLProbability{0.1},\qquad
p_X(\GTenHLValue{1})=\GTenHLProbability{0.5},\qquad
p_X(\GTenHLValue{2})=\GTenHLProbability{0.4}.
\]
\[
\sum_{x\in S}p_X(x)=p_X(\GTenHLValue{0})+p_X(\GTenHLValue{1})+p_X(\GTenHLValue{2}).
\]
This sum runs over numerical values. The earlier sum $\sum_{\omega:X(\omega)=x}P(\omega)$ runs over the outcomes that produce one chosen value. Group those outcomes first, then include each value's probability once. In particular, include $p_X(\GTenHLValue{1})=\GTenHLProbability{0.5}$ once, even though both $b$ and $c$ produce 1.
\item \textbf{Read the final $=1$ and check it.} The supported values are mutually exclusive and together cover all possible values of $X$, so their probabilities must add to 1:
\[
p_X(\GTenHLValue{0})+p_X(\GTenHLValue{1})+p_X(\GTenHLValue{2})=\GTenHLProbability{0.1}+\GTenHLProbability{0.5}+\GTenHLProbability{0.4}=1.
\]
The total 1 means a 100\% chance of obtaining some supported value. Financially, the fee must be \GTenHLValue{\$0}, \GTenHLValue{\$1}, or \GTenHLValue{\$2}; the total probability is 1, not a payment of \$1. Checking this total is called \emph{normalization}.
\end{enumerate}

For a finite support, the sum has a last term. For a countably infinite support, such as $S=\{1,2,3,\ldots\}$, the same notation means adding an infinite list:
\[
\sum_{x\in S}p_X(x)=p_X(1)+p_X(2)+p_X(3)+\cdots=1.
\]
More precisely, the totals of the first finitely many terms must approach 1 as more terms are included. Every term must still be non-negative. An infinite list of possible values can therefore have a finite total probability of 1; the next problems show how to check both finite and infinite models.
\Needspace{8\baselineskip}
\subsection{Advanced general problem: a finite model}
\subsubsection*{Mathematical wording}
A model assigns masses $\GTenHLParameter{c},\GTenHLParameter{2c},\GTenHLParameter{3c}$ to scores $\GTenHLValue{0,2,5}$, respectively; let $\GTenHLValue{X}$ be its score.

\textbf{Questions/Tasks}
\begin{enumerate}
\item Determine $\GTenHLParameter{c}$, verify the PMF, and calculate $P(X\geq2)$.
\end{enumerate}
\Needspace{6\baselineskip}
\subsubsection*{Financial narrative wording}
A shop offers a promotional refund on one purchase. The refund can only be \GTenHLValue{\$0}, \GTenHLValue{\$2}, or \GTenHLValue{\$5}. A \$2 refund is twice as likely as no refund, and a \$5 refund is three times as likely as no refund. These are the only possible refund amounts.

\textbf{Questions/Tasks}
\begin{enumerate}
\item Find the chance of no refund and use it to find the chances of the other refund amounts. Check that all refund chances are non-negative and add to a total chance of 1. What is the chance of a refund of at least \$2?
\end{enumerate}
\subsubsection*{Complete worked solution}
\par\noindent\begin{minipage}{\linewidth}
\textbf{Step 1: Define and identify support.} $\GTenHLValue{X}$ is the score with support $\GTenHLValue{\{0,2,5\}}$. In the financial story, $\GTenHLValue{X}$ is the \GTenHLValue{refund in USD}. If $\GTenHLParameter{c}$ is the chance of no refund, the relative likelihoods give masses $\GTenHLParameter{c},\GTenHLParameter{2c},\GTenHLParameter{3c}$ for refunds of 0, 2, and 5 USD. This is the same finite model; normalization must be checked.
\end{minipage}\par\medskip
\par\noindent\begin{minipage}{\linewidth}
\textbf{Step 2: Normalize.} $\sum_xp_X(x)=\GTenHLParameter{c}+\GTenHLParameter{2c}+\GTenHLParameter{3c}=\GTenHLParameter{6c}=1$, so $\GTenHLParameter{c}=\GTenHLParameter{1/6}$. The masses are $\GTenHLProbability{1/6},\GTenHLProbability{1/3},\GTenHLProbability{1/2}$, all positive and with total 1, so they define a PMF. In the story, these are the chances of refunds of \$0, \$2, and \$5, respectively.
\end{minipage}\par\medskip
\par\noindent\begin{minipage}{\linewidth}
\textbf{Step 3: Translate and calculate.} $P(\GTenHLValue{X}\geq\GTenHLValue{2})=p_X(\GTenHLValue{2})+p_X(\GTenHLValue{5})=\GTenHLProbability{1/3}+\GTenHLProbability{1/2}=\GTenHLProbability{5/6}\approx\GTenHLProbability{0.8333}$.
\end{minipage}\par\medskip
\par\noindent\begin{minipage}{\linewidth}
\textbf{Step 4: Check and interpret.} The probability lies between 0 and 1. The model has only three possible values despite the gap from 2 to 5. Financially, the chance of a refund of at least \$2 is approximately \GTenHLProbability{83.33\%}. The refund is measured in USD; the probability has no units.
\end{minipage}\par\medskip
% Graphs follow completed calculations; captions stay with their panels.
\par\noindent \begin{minipage}[t]{.485\linewidth}
\centering
\GTenPMFGraph[width=\linewidth,height=4.7cm,xlabel={Score / refund (USD)}]{0,2,5}{1/6,1/3,1/2}
\par\small \textbf{Normalized finite PMF.} Only 0, 2, and 5 are supported; the gap is preserved. Heights display rounded probabilities.
\end{minipage}\par\medskip

\Needspace{8\baselineskip}
\subsection{Advanced general problem: an infinite model}
\subsubsection*{Mathematical wording}
A model assigns $p_T(t)=2^{-t}$ for $t=1,2,\ldots$, where $\GTenHLValue{T}$ is a \GTenHLValue{positive-integer trial count}.

\textbf{Questions/Tasks}
\begin{enumerate}
\item Verify the PMF and calculate $P(T>3)$.
\end{enumerate}
\Needspace{8\baselineskip}
\subsubsection*{Financial narrative wording}
A student running a school enterprise calls potential customers until the first sale. Each call is an independent attempt, with a constant one-half chance of making a sale. Count every call, including the call that makes the first sale. There is no limit on the number of calls allowed, and calling stops at the first sale.

\textbf{Questions/Tasks}
\begin{enumerate}
\item Check that the chances for every possible number of calls up to and including the first sale are non-negative and add to 1, allowing for arbitrarily many calls. What is the chance that \GTenHLValue{more than 3 calls} are needed to make the first sale?
\end{enumerate}
\subsubsection*{Complete worked solution}
\par\noindent\begin{minipage}{\linewidth}
\textbf{Step 1: Define and identify support.} $\GTenHLValue{T}$ is a \GTenHLValue{positive-integer trial count} with support $\GTenHLValue{\{1,2,\ldots\}}$. In the financial story, $\GTenHLValue{T}$ counts calls, including the successful call. A first sale on call $\GTenHLValue{t}$ requires $t-1$ failures followed by a success. Independence gives $p_T(t)=(1/2)^{t-1}(1/2)=2^{-t}$. This is the same infinite model; normalization must be checked.
\end{minipage}\par\medskip
\par\noindent\begin{minipage}{\linewidth}
\textbf{Step 2: Normalize.} The first $m$ terms total $1-2^{-m}$; as $m$ grows without bound this approaches 1. Every $2^{-t}$ is positive, so the infinite list defines a PMF. In the story, the total chance of a first sale on some finite call is 1.
\end{minipage}\par\medskip
\par\noindent\begin{minipage}{\linewidth}
\textbf{Step 3: Translate and calculate.} $P(\GTenHLValue{T}>\GTenHLValue{3})=1-\sum_{t=1}^{3}\GTenHLProbability{2^{-t}}=1-(\GTenHLProbability{1/2}+\GTenHLProbability{1/4}+\GTenHLProbability{1/8})=\GTenHLProbability{1/8}=\GTenHLProbability{0.1250}$.
\end{minipage}\par\medskip
\par\noindent\begin{minipage}{\linewidth}
\textbf{Step 4: Check and interpret.} The probability lies between 0 and 1. There is no last possible trial; infinite support does not mean infinite total probability. Financially, the chance of needing \GTenHLValue{more than 3 calls} for the first sale is \GTenHLProbability{12.5\%}. The waiting count is measured in calls; the probability has no units.
\end{minipage}\par\medskip
% Graphs follow completed calculations; captions stay with their panels.
\par\noindent \begin{minipage}[t]{.98\linewidth}
\centering
\GTenPMFGraph[width=\linewidth,height=4.7cm,xlabel={$t$ (calls), continuing to $9,10,\ldots$},tail=1/256]{1,2,3,4,5,6,7,8}{1/2,1/4,1/8,1/16,1/32,1/64,1/128,1/256}
\par\small \textbf{Waiting-count PMF: first eight values only.} $p_T(t)=2^{-t}$ for $t=1,2,\ldots$; the support continues indefinitely to the right. The omitted probability is $P(T>8)=\sum_{t=9}^{\infty}2^{-t}=1/256$. Displayed bars total $255/256$ and are not renormalized; labels are rounded.
\end{minipage}\par\medskip

\section{Notation reference and practice procedure}
\begin{center}
\begin{tabular}{lp{0.72\linewidth}}
\toprule
Symbol & Meaning and units \\ \midrule
$\omega$ & One complete experimental outcome \\
$X,x$ & Random variable and one numerical value \\
$S$ & Support: values with positive probability \\
$p_X(x)$ & Probability assigned to value $x$ \\
$\sum$ & Add the indicated terms \\
\bottomrule
\end{tabular}
\end{center}
\begin{enumerate}
\item Define the random variable in words and give its units and support.
\item List parameters and check that the model assumptions match the situation.
\item Translate ordinary language into an event, paying attention to endpoints.
\item Write the symbolic formula first, then substitute without rounding intermediate values.
\item Check bounds, total probability or a complementary calculation; interpret the answer in context.
\end{enumerate}
\Needspace{10\baselineskip}
\section{Foundational contextual practice}\label{sec:t2-c1-foundational}
\Needspace{8\baselineskip}
\subsection{Notebook buyers}
Two visitors buy independently with probability $1/2$ each. Let B mean a purchase and N no purchase. Each sale earns 4 USD.

\textbf{Questions/Tasks}
\begin{enumerate}
\item List outcomes, define the buyer count $\GTenHLValue{X}$, build its PMF and find the chance of revenue 4 USD.
\end{enumerate}
\subsubsection*{Complete worked solution}
\par\noindent\begin{minipage}{\linewidth}
\textbf{Step 1: Define and list.} $\GTenHLValue{X}$ counts buyers, with support $\GTenHLValue{\{0,1,2\}}$. Revenue is $\GTenHLSecond{Y}=4\GTenHLValue{X}$ USD, with support $\GTenHLSecond{\{0,4,8\}}$ USD.
\end{minipage}\par\medskip
\par\noindent\begin{minipage}{\linewidth}
\textbf{Step 2: Check the model and map.} Independence and the two equal probabilities make BB, BN, NB, NN each have probability $(1/2)^2=1/4$. Their $\GTenHLValue{X}$ values are 2, 1, 1, 0.
\end{minipage}\par\medskip
\par\noindent\begin{minipage}{\linewidth}
\textbf{Step 3: Calculate.} $p_X(\GTenHLValue{1})=P(\GTenHLOutcome{\mathrm{BN}})+P(\GTenHLOutcome{\mathrm{NB}})=\GTenHLProbability{1/4}+\GTenHLProbability{1/4}=\GTenHLProbability{0.5}$.\begin{center}
\begin{tabular}{llll}
\toprule
$\GTenHLValue{x}$ (buyers) & \GTenHLValue{0} & \GTenHLValue{1} & \GTenHLValue{2} \\ \midrule
$p_X(x)$ & \GTenHLProbability{0.25} & \GTenHLProbability{0.50} & \GTenHLProbability{0.25} \\
\bottomrule
\end{tabular}
\end{center}
$P(\GTenHLSecond{Y}=\GTenHLSecond{4})=P(\GTenHLValue{X}=\GTenHLValue{1})=\GTenHLProbability{0.5000}$.
\end{minipage}\par\medskip
\par\noindent\begin{minipage}{\linewidth}
\textbf{Step 4: Check and interpret.} The table totals 1. A revenue of \GTenHLSecond{4 USD} occurs in two of the four equally likely outcomes; 0.5 is a probability, not revenue.
\end{minipage}\par\medskip
% Graphs follow completed calculations; captions stay with their panels.
\par\noindent \begin{minipage}[t]{.485\linewidth}
\centering
\GTenOutcomeProbabilityGraph[width=\linewidth,height=4.7cm]{BB,BN,NB,NN}{.25,.25,.25,.25}
\par\small \textbf{Complete-outcome probabilities.} Independence makes all four outcomes equally likely.
\end{minipage}\hfill\begin{minipage}[t]{.485\linewidth}
\centering
\GTenOutcomeValueGraph[width=\linewidth,height=4.7cm,ylabel={$X(\omega)$ (buyers)}]{BB,BN,NB,NN}{2,1,1,0}
\par\small \textbf{Assigned buyer counts.} The bars show $X(\omega)$, measured in buyers.
\end{minipage}\par\medskip
\par\noindent \begin{minipage}[t]{.485\linewidth}
\centering
\GTenPMFGraph[width=\linewidth,height=4.7cm,xlabel={$x$ (buyers)}]{0,1,2}{.25,.5,.25}
\par\small \textbf{Buyer-count PMF.} BN and NB are grouped at one buyer.
\end{minipage}\hfill\begin{minipage}[t]{.485\linewidth}
\centering
\GTenPMFGraph[width=\linewidth,height=4.7cm,xlabel={$y$ (USD)},xstep=4]{0,4,8}{.25,.5,.25}
\par\small \textbf{Revenue PMF.} $\GTenHLSecond{Y}=4\GTenHLValue{X}$ in USD; the corresponding probabilities are unchanged.
\end{minipage}\par\medskip

\secInicio[][sec:t2-c1-foundational]
\Needspace{8\baselineskip}
\subsection{Daily insurance claims}
The number $\GTenHLValue{X}$ of claims on a small policy group in one day has probabilities $\GTenHLProbability{0.7},\GTenHLParameter{q},\GTenHLProbability{0.1}$ for $\GTenHLValue{0,1,2}$ claims; no other values are possible.

\textbf{Questions/Tasks}
\begin{enumerate}
\item Find $\GTenHLParameter{q}$, verify the PMF and calculate the probability of any claim.
\end{enumerate}
\subsubsection*{Complete worked solution}
\par\noindent\begin{minipage}{\linewidth}
\textbf{Step 1: Define.} $\GTenHLValue{X}$ counts \GTenHLValue{daily claims}, with proposed support $\GTenHLValue{\{0,1,2\}}$. The supplied model is a complete finite distribution; no trial independence assumption is required to use its table.
\end{minipage}\par\medskip
\par\noindent\begin{minipage}{\linewidth}
\textbf{Step 2: Normalize.} $\GTenHLParameter{q}=1-p_X(\GTenHLValue{0})-p_X(\GTenHLValue{2})=1-\GTenHLProbability{0.7}-\GTenHLProbability{0.1}=\GTenHLParameter{0.2}$.
\end{minipage}\par\medskip
\par\noindent\begin{minipage}{\linewidth}
\textbf{Step 3: Translate and calculate.} $P(\GTenHLValue{X}\geq\GTenHLValue{1})=p_X(\GTenHLValue{1})+p_X(\GTenHLValue{2})=\GTenHLProbability{0.2}+\GTenHLProbability{0.1}=\GTenHLProbability{0.3000}$.
\end{minipage}\par\medskip
\par\noindent\begin{minipage}{\linewidth}
\textbf{Step 4: Check and interpret.} All masses are positive and total 1. Any claim occurs with probability \GTenHLProbability{30\%}; the count itself is measured in claims, not USD.
\end{minipage}\par\medskip
% Graphs follow completed calculations; captions stay with their panels.
\par\noindent \begin{minipage}[t]{.485\linewidth}
\centering
\GTenPMFGraph[width=\linewidth,height=4.7cm,xlabel={$x$ (claims)}]{0,1,2}{.7,.2,.1}
\par\small \textbf{Completed claim-count PMF.} The missing mass is $q=0.2$.
\end{minipage}\par\medskip

\secInicio[][sec:t2-c1-foundational]
\Needspace{10\baselineskip}
\section{Intermediate contextual practice}\label{sec:t2-c1-intermediate}
\Needspace{8\baselineskip}
\subsection{Payment methods and a fee}
A shop receives one payment: \GTenHLOutcome{cash} with probability \GTenHLProbability{0.25}, \GTenHLOutcome{debit} with \GTenHLProbability{0.45}, or \GTenHLOutcome{credit} with \GTenHLProbability{0.30}. Cash costs no processing fee; either card costs 1 USD.

\textbf{Questions/Tasks}
\begin{enumerate}
\item Define the fee variable and construct its PMF. Explain why three outcomes need not give three values.
\end{enumerate}
\subsubsection*{Complete worked solution}
\par\noindent\begin{minipage}{\linewidth}
\textbf{Step 1: Define.} $\GTenHLValue{X}$ is the \GTenHLValue{processing fee in USD}, with support $\GTenHLValue{\{0,1\}}$. The method categories are mutually exclusive and exhaustive.
\end{minipage}\par\medskip
\par\noindent\begin{minipage}{\linewidth}
\textbf{Step 2: Map.} $\GTenHLValue{X}(\GTenHLOutcome{\text{cash}})=\GTenHLValue{0}$ and $\GTenHLValue{X}(\GTenHLOutcome{\text{debit}})=\GTenHLValue{X}(\GTenHLOutcome{\text{credit}})=\GTenHLValue{1}$. Add the probabilities of the two card outcomes.
\end{minipage}\par\medskip
\par\noindent\begin{minipage}{\linewidth}
\textbf{Step 3: Calculate.} $p_X(\GTenHLValue{0})=\GTenHLProbability{0.25}$ and $p_X(\GTenHLValue{1})=P(\GTenHLOutcome{\text{debit}})+P(\GTenHLOutcome{\text{credit}})=\GTenHLProbability{0.45}+\GTenHLProbability{0.30}=\GTenHLProbability{0.75}$. Thus $P(\GTenHLValue{X}>\GTenHLValue{0})=\GTenHLProbability{0.7500}$.
\end{minipage}\par\medskip
\par\noindent\begin{minipage}{\linewidth}
\textbf{Step 4: Check and interpret.} $0.25+0.75=1$. There are three method outcomes but only two fee values; a random variable need not assign different values to different outcomes.
\end{minipage}\par\medskip
% Graphs follow completed calculations; captions stay with their panels.
\par\noindent \begin{minipage}[t]{.32\linewidth}
\centering
\GTenOutcomeProbabilityGraph[width=\linewidth,height=4.7cm]{cash,debit,credit}{.25,.45,.30}
\par\small \textbf{Payment-method probabilities.} The bars refer to method outcomes.
\end{minipage}\hfill\begin{minipage}[t]{.32\linewidth}
\centering
\GTenOutcomeValueGraph[width=\linewidth,height=4.7cm,ylabel={$X(\omega)$ (USD)}]{cash,debit,credit}{0,1,1}
\par\small \textbf{Assigned fees.} Both cards map to 1 USD; cash maps to zero.
\end{minipage}\hfill\begin{minipage}[t]{.32\linewidth}
\centering
\GTenPMFGraph[width=\linewidth,height=4.7cm,xlabel={$x$ (USD)}]{0,1}{.25,.75}
\par\small \textbf{Grouped fee PMF.} Debit and credit contribute $0.45+0.30=0.75$ at 1 USD.
\end{minipage}\par\medskip


\secInicio[][sec:t2-c1-intermediate]
\Needspace{8\baselineskip}
\subsection{Investment results with a missing mass}
A hypothetical one-month net gain $\GTenHLValue{G}$ takes values $\GTenHLValue{-20,0,30}$ USD with masses $\GTenHLProbability{0.2},\GTenHLProbability{0.5},\GTenHLParameter{q}$. A student proposes $\GTenHLParameter{q}=\GTenHLParameter{0.4}$.

\textbf{Questions/Tasks}
\begin{enumerate}
\item Correct the PMF, identify the support and calculate the probability of a non-negative gain.
\end{enumerate}
\subsubsection*{Complete worked solution}
\par\noindent\begin{minipage}{\linewidth}
\textbf{Step 1: Define.} $\GTenHLValue{G}$ is monetary net gain, not a count. A finite set of monetary outcomes is discrete even though money is involved.
\end{minipage}\par\medskip
\par\noindent\begin{minipage}{\linewidth}
\textbf{Step 2: Validate.} $\GTenHLParameter{q}=1-\GTenHLProbability{0.2}-\GTenHLProbability{0.5}=\GTenHLParameter{0.3}$; the proposed 0.4 would make total mass 1.1, so it is invalid. Support is $\GTenHLValue{\{-20,0,30\}}$ USD.
\end{minipage}\par\medskip
\par\noindent\begin{minipage}{\linewidth}
\textbf{Step 3: Calculate.} $P(\GTenHLValue{G}\geq\GTenHLValue{0})=p_G(\GTenHLValue{0})+p_G(\GTenHLValue{30})=\GTenHLProbability{0.5}+\GTenHLProbability{0.3}=\GTenHLProbability{0.8000}$.
\end{minipage}\par\medskip
\par\noindent\begin{minipage}{\linewidth}
\textbf{Step 4: Check and interpret.} The corrected masses are non-negative and total 1. A loss remains possible with probability \GTenHLProbability{0.2}. Zero gain is included in non-negative gain.
\end{minipage}\par\medskip
% Graphs follow completed calculations; captions stay with their panels.
\par\noindent \begin{minipage}[t]{.485\linewidth}
\centering
\GTenPMFGraph[width=\linewidth,height=4.7cm,xlabel={$g$ (USD)},xstep=10]{-20,0,30}{.2,.5,.3}
\par\small \textbf{Corrected gain PMF.} $q=0.3$; negative, zero, and positive gains retain their actual numerical spacing.
\end{minipage}\par\medskip

\secInicio[][sec:t2-c1-intermediate]
\Needspace{10\baselineskip}
\section{Advanced contextual practice}\label{sec:t2-c1-advanced}
\Needspace{8\baselineskip}
\subsection{Calls until a sale}
\textit{Markers here:} \GTenHLValue{$T$, call count}; \GTenHLSecond{$D$, duration in minutes}.
A sales team models $\GTenHLValue{T}$, the call number of its first sale, by $p_T(t)=(1/2)^t$ for $t=1,2,\ldots$. It also records \GTenHLSecond{exact call duration} $\GTenHLSecond{D}$ in minutes.

\textbf{Questions/Tasks}
\begin{enumerate}
\item Classify both variables; verify the waiting-count PMF and find $P(T\leq3)$.
\end{enumerate}
\subsubsection*{Complete worked solution}
\par\noindent\begin{minipage}{\linewidth}
\textbf{Step 1: Define and distinguish.} $\GTenHLValue{T}$ counts calls and has countably infinite support. $\GTenHLSecond{D}$ measures time and is treated as continuous; its exact support contains intervals, rather than a list of isolated counts.
\end{minipage}\par\medskip
\par\noindent\begin{minipage}{\linewidth}
\textbf{Step 2: Check the supplied model.} All masses $2^{-t}$ are positive. Their first $m$ terms sum to $1-2^{-m}\to1$, so the infinite PMF is valid. It is compatible with independent calls with constant sale chance $1/2$.
\end{minipage}\par\medskip
\par\noindent\begin{minipage}{\linewidth}
\textbf{Step 3: Calculate.} $P(\GTenHLValue{T}\leq\GTenHLValue{3})=\sum_{t=1}^{3}p_T(t)=\GTenHLProbability{1/2}+\GTenHLProbability{1/4}+\GTenHLProbability{1/8}=\GTenHLProbability{7/8}=\GTenHLProbability{0.8750}$.
\end{minipage}\par\medskip
\par\noindent\begin{minipage}{\linewidth}
\textbf{Step 4: Interpret.} The first sale occurs \GTenHLValue{within three calls} with probability \GTenHLProbability{87.5\%}. There is no fixed maximum call number in the model; this does not make the count continuous.
\end{minipage}\par\medskip
% Graphs follow completed calculations; captions stay with their panels.
\par\noindent \begin{minipage}[t]{.98\linewidth}
\centering
\GTenPMFGraph[width=\linewidth,height=4.7cm,xlabel={$t$ (calls), continuing to $9,10,\ldots$},tail=1/256]{1,2,3,4,5,6,7,8}{1/2,1/4,1/8,1/16,1/32,1/64,1/128,1/256}
\par\small \textbf{Waiting-count PMF: first eight values only.} $p_T(t)=2^{-t}$ for $t=1,2,\ldots$; the support continues indefinitely to the right. The omitted probability is $P(T>8)=\sum_{t=9}^{\infty}2^{-t}=1/256$. Displayed bars total $255/256$ and are not renormalized; labels are rounded.
\end{minipage}\par\medskip

\secInicio[][sec:t2-c1-advanced]
\Needspace{8\baselineskip}
\subsection{Delivery records: two numerical descriptions}
\textit{Markers here:} \GTenHLValue{$B$, bonus in USD}; \GTenHLSecond{$I$, lateness indicator}.
A delivery has outcomes \GTenHLOutcome{early}, \GTenHLOutcome{on time}, \GTenHLOutcome{late} with probabilities $\GTenHLProbability{0.2},\GTenHLProbability{0.5},\GTenHLProbability{0.3}$. A bonus $\GTenHLValue{B}$ is 2 USD for early or on time and $-3$ USD for late. An indicator $\GTenHLSecond{I}$ is 1 if late and 0 otherwise.

\textbf{Questions/Tasks}
\begin{enumerate}
\item Build both PMFs and contrast them with an exact delivery-time measurement.
\end{enumerate}
\subsubsection*{Complete worked solution}
\par\noindent\begin{minipage}{\linewidth}
\textbf{Step 1: Define and map.} $\GTenHLValue{B}$ measures \GTenHLValue{bonus in USD}, with support $\GTenHLValue{\{-3,2\}}$; $\GTenHLSecond{I}$ labels lateness, with support $\GTenHLSecond{\{0,1\}}$. Early and on time both map to $\GTenHLValue{B}=\GTenHLValue{2}$ and $\GTenHLSecond{I}=\GTenHLSecond{0}$.
\end{minipage}\par\medskip
\par\noindent\begin{minipage}{\linewidth}
\textbf{Step 2: Check the categories.} The three exclusive outcomes cover all deliveries and their probabilities total 1. No equal-likelihood assumption is justified.
\end{minipage}\par\medskip
\par\noindent\begin{minipage}{\linewidth}
\textbf{Step 3: Calculate.} $p_B(\GTenHLValue{2})=\GTenHLProbability{0.2}+\GTenHLProbability{0.5}=\GTenHLProbability{0.7}$, $p_B(\GTenHLValue{-3})=\GTenHLProbability{0.3}$; $p_I(\GTenHLSecond{0})=\GTenHLProbability{0.7}$, $p_I(\GTenHLSecond{1})=\GTenHLProbability{0.3}$. Thus $P(\GTenHLValue{B}<\GTenHLValue{0})=P(\GTenHLSecond{I}=\GTenHLSecond{1})=\GTenHLProbability{0.3000}$.
\end{minipage}\par\medskip
\par\noindent\begin{minipage}{\linewidth}
\textbf{Step 4: Check and interpret.} Each PMF totals 1. The same uncertain delivery can generate different discrete variables with different units. Exact elapsed delivery time is a continuous measurement under an ideal model, even though its early/on-time/late classification is discrete.
\end{minipage}\par\medskip
% Graphs follow completed calculations; captions stay with their panels.
\par\noindent \begin{minipage}[t]{.32\linewidth}
\centering
\GTenOutcomeProbabilityGraph[width=\linewidth,height=4.1cm]{early,on time,late}{.2,.5,.3}
\par\small \textbf{Delivery-outcome probabilities.} The bars refer to delivery categories.
\end{minipage}\hfill\begin{minipage}[t]{.32\linewidth}
\centering
\GTenOutcomeValueGraph[width=\linewidth,height=4.1cm,ylabel={$B(\omega)$ (USD)}]{early,on time,late}{2,2,-3}
\par\small \textbf{Assigned bonuses.} Late delivery gives $-3$ USD, shown below zero.
\end{minipage}\hfill\begin{minipage}[t]{.32\linewidth}
\centering
\GTenOutcomeValueGraph[width=\linewidth,height=4.1cm,ylabel={$I(\omega)$ (indicator)}]{early,on time,late}{0,0,1}
\par\small \textbf{Assigned lateness indicator.} Only late delivery maps to 1; the indicator has no units.
\end{minipage}\par\medskip
\par\noindent \begin{minipage}[t]{.485\linewidth}
\centering
\GTenPMFGraph[width=\linewidth,height=4.1cm,xlabel={$b$ (USD)}]{-3,2}{.3,.7}
\par\small \textbf{Bonus PMF.} Early and on time group at 2 USD; a negative bonus still has a positive probability.
\end{minipage}\hfill\begin{minipage}[t]{.485\linewidth}
\centering
\GTenPMFGraph[width=\linewidth,height=4.1cm,xlabel={$i$ (indicator)}]{0,1}{.7,.3}
\par\small \textbf{Indicator PMF.} Early and on time group at 0. Compare the outcome probabilities above.
\end{minipage}\par\medskip


\secInicio[][sec:t2-c1-advanced]
\section{Closing reference: key ideas and common errors}
An outcome is a whole story; a random variable maps it to a number. Group outcomes before assigning masses. Support contains only positive-mass values. Discrete does not mean equally likely, and a monetary variable need not be continuous. Check every mass and the total; never confuse a number of outcomes with a probability.
\secInicio
\end{document}

